\section{Дальность связи}

Дальность определяется тем, дойдёт ли до приёмника сигнал, который \textbf{сильнее его порога чувствительности} (и сильнее помех). Это считают через \term{бюджет радиолинии}.

\subsection{Децибелы — быстрая памятка}
\begin{itemize}
\item \term{дБм} — мощность относительно \SI{1}{\milli\watt}: $P_\text{дБм} = 10\lg\frac{P}{1\,\text{мВт}}$. \SI{25}{\milli\watt} = 14~дБм; \SI{100}{\milli\watt} = 20~дБм; \SI{250}{\milli\watt} = 24~дБм; \SI{1}{\watt} = 30~дБм.
\item \textbf{+3~дБ} = мощность ×2; \textbf{+10~дБ} = ×10; \textbf{+6~дБ} = дальность ×2 (в свободном пространстве).
\item \term{дБи} — усиление антенны относительно идеального изотропного излучателя. Диполь — 2~дБи, патч — 8–10~дБи, helical — 10–14~дБи.
\end{itemize}

\subsection{Потери в свободном пространстве и бюджет линии}
\begin{formula}[Потери в свободном пространстве (FSPL)]
\[
L_\text{FSPL}\,[\text{дБ}] = 20\lg d_\text{км} + 20\lg f_\text{МГц} + 32{,}44
\]
Потери растут с \textbf{расстоянием} и с \textbf{частотой}: удвоение расстояния даёт +6~дБ потерь. При одинаковой дальности на 900~МГц потери на $\approx$\,8,5~дБ меньше, чем на 2,4~ГГц.
\end{formula}

\begin{formula}[Бюджет радиолинии]
\[
P_\text{пр} = P_\text{пер} + G_\text{пер} + G_\text{пр} - L_\text{FSPL} - L_\text{доп} \;\geq\; S_\text{пр} + M
\]
$P_\text{пер}$ — мощность передатчика (дБм), $G$ — усиления антенн (дБи), $L_\text{доп}$ — потери в кабелях, на препятствиях, из-за поляризации, $S_\text{пр}$ — \textbf{чувствительность приёмника} (дБм, например $-105$), $M$ — запас на замирания (10–20~дБ).
\end{formula}

\begin{figure}[H]\centering
\begin{tikzpicture}
\begin{semilogxaxis}[width=0.85\linewidth,height=62mm,xlabel={Расстояние, км},ylabel={Потери FSPL, дБ},
  xmin=0.1,xmax=100,ymin=65,ymax=150,grid=both,grid style={gray!20},legend pos=south east,legend style={font=\small\sffamily},
  xticklabels={0{,}1,1,10,100},xtick={0.1,1,10,100}]
\addplot[very thick,okgreen,domain=0.1:100,samples=50] {20*log10(x)+20*log10(868)+32.44};
\addplot[very thick,accent,domain=0.1:100,samples=50] {20*log10(x)+20*log10(2400)+32.44};
\addplot[very thick,bad,domain=0.1:100,samples=50] {20*log10(x)+20*log10(5800)+32.44};
\legend{{868 МГц},{2,4 ГГц},{5,8 ГГц}}
\end{semilogxaxis}
\end{tikzpicture}
\caption{Потери в свободном пространстве для частот RC-линков и видео}
\end{figure}

\subsection{Пример расчёта}
\textbf{ELRS 2,4~ГГц, 250~мВт, 500~Гц.} $P_\text{пер} = 24$~дБм, антенны $2+2$~дБи, чувствительность на 500~Гц $\approx -105$~дБм, запас 10~дБ.
\[
L_\text{max} = 24 + 2 + 2 - (-105) - 10 = 123~\text{дБ}
\;\Rightarrow\;
\lg d = \frac{123 - 32{,}44 - 67{,}6}{20} \approx 1{,}15
\;\Rightarrow\; d \approx 14~\text{км}.
\]
Если перейти на 50~Гц (чувствительность $\approx -115$~дБм), допустимые потери +10~дБ, и теоретическая дальность вырастает в $10^{10/20}\approx3{,}2$ раза — до $\sim$\,45~км. Это \emph{теория для прямой видимости}: на практике мешают земля, препятствия, помехи, ориентация антенн.

\textbf{Аналоговый VTX 5,8~ГГц, 600~мВт.} $P_\text{пер} \approx 28$~дБм, антенны 2~дБи (на дроне) + 8~дБи (патч), чувствительность видеоприёмника $\approx -90$~дБм, запас 10~дБ: $L_\text{max} = 118$~дБ; на 5,8~ГГц это $d \approx 3$~км.

\subsection{Зона Френеля и радиогоризонт}
\begin{figure}[H]\centering
\begin{tikzpicture}[font=\sffamily\small]
\fill[brown!25] (-0.5,0) rectangle (10.5,-0.4);
\draw[thick] (-0.5,0) -- (10.5,0);
\draw[thick] (0,0) -- (0,1.5) node[above] {пилот};
\draw[thick] (10,0) -- (10,2.2) node[above] {дрон};
\draw[fill=accent!15,draw=accent,thick] (0,1.5) .. controls (3,3.3) and (7,3.9) .. (10,2.2) .. controls (7,0.5) and (3,-0.1) .. (0,1.5);
\draw[dashed,thick] (0,1.5) -- (10,2.2);
\fill[okgreen!60!black] (4.2,0) -- (4.8,1.2) -- (5.4,0) -- cycle;
\node[text=bad,font=\sffamily\scriptsize,align=center] at (6.4,0.6) {препятствие\\в зоне Френеля};
\draw[<->] (5,1.85) -- node[right,font=\scriptsize] {$r$} (5,2.9);
\end{tikzpicture}
\caption{Первая зона Френеля должна быть свободна хотя бы на 60\,\%}
\end{figure}

\begin{formula}[Радиус первой зоны Френеля (в середине трассы) и радиогоризонт]
\[
r\,[\text{м}] \approx 8{,}66\sqrt{\frac{D_\text{км}}{f_\text{ГГц}}}
\qquad\qquad
d_\text{гор}\,[\text{км}] \approx 4{,}12\left(\sqrt{h_1} + \sqrt{h_2}\right),\; h\text{ — в метрах}
\]
Пример: 2,4~ГГц, 5~км $\Rightarrow r \approx 12{,}5$~м. Поэтому даже при «прямой видимости» низко летящий дрон теряет сигнал: земля и деревья заходят в зону Френеля. Для 900~МГц $r$ ещё больше ($\approx 20$~м).\\
Горизонт: пилот $h_1 = 2$~м, дрон $h_2 = 100$~м $\Rightarrow d \approx 4{,}12\,(1{,}4 + 10) \approx 47$~км.
\end{formula}

\subsection{От чего зависит реальная дальность}
\begin{enumerate}
\item \textbf{Мощность передатчика} — +6~дБ (×4 мощности) = ×2 дальности.
\item \textbf{Чувствительность приёмника} — зависит от модуляции и скорости: LoRa на низкой частоте пакетов чувствительнее всего.
\item \textbf{Частота} — ниже частота $\Rightarrow$ меньше потери и лучше огибание препятствий.
\item \textbf{Антенны}: усиление, \textbf{совпадение поляризации} (линейная–круговая: $-3$~дБ; RHCP–LHCP: $-20$~дБ и больше), ориентация (у диполя «мёртвая зона» вдоль оси), расположение (не закрывать карбоном и батареей).
\item \textbf{Высота полёта и препятствия} — зона Френеля, радиогоризонт, здания, лес, рельеф.
\item \textbf{Помехи}: Wi-Fi, другие пилоты, городской эфир, ЛЭП, \textbf{средства РЭБ}.
\item \textbf{Самое слабое звено}: дальность ограничивает тот канал, который отказывает первым — обычно \textbf{видео}, а не управление. И ещё — ёмкость аккумулятора: нужно успеть вернуться.
\end{enumerate}

\begin{table}[H]\centering\small
\renewcommand{\arraystretch}{1.2}
\begin{tabular}{@{}lr@{}}
\toprule
\textbf{Система} & \textbf{Ориентировочная реальная дальность} \\
\midrule
Wi-Fi-дрон (игрушка) & \SIrange{50}{300}{\metre} \\
Аналоговый VTX 25 мВт, всенаправленные антенны & \SIrange{0.3}{1}{\kilo\metre} \\
Аналоговый VTX 600–1000 мВт, патч/helical на земле & \SIrange{3}{8}{\kilo\metre} \\
Цифровой DJI O3/O4 & до \SIrange{10}{15}{\kilo\metre} \\
FrSky ACCST 2,4 ГГц, 100 мВт & \SIrange{1.5}{2}{\kilo\metre} \\
ELRS 2,4 ГГц 250 мВт (250–500 Гц) & \SIrange{5}{15}{\kilo\metre} \\
ELRS 900 / Crossfire, 250 мВт–1 Вт & \SIrange{20}{50}{\kilo\metre}+ \\
\bottomrule
\end{tabular}
\caption{Типичная дальность (в хороших условиях, без помех)}
\end{table}

\begin{keybox}[Как увеличить дальность]
Поднять антенну наземной станции повыше; использовать направленную антенну (и трекер); правильно ориентировать и разнести антенны на дроне; понизить частоту пакетов (ELRS 50~Гц); перейти на 868~МГц; летать выше (зона Френеля, горизонт); использовать ретранслятор; следить за LQ и RSSI в OSD и вовремя поворачивать назад.
\end{keybox}

\begin{exam}
\begin{enumerate}
\item Во сколько раз изменится дальность, если мощность передатчика увеличить в 4 раза?
\item Посчитайте потери FSPL на 2,4~ГГц на расстоянии 2~км.
\item Что такое зона Френеля и почему низко летящий дрон теряет связь?
\item Почему линейная антенна на передатчике и круговая на приёмнике — плохая пара?
\end{enumerate}
\end{exam}
